% 概念白话解释（零基础导读）。由 main.tex \input，不单独编译。
一句话搭骨架：Polymarket 是一个\textbf{押注"未来某件事会不会发生"}的
网站；本研究把它上面所有成交记录（谁、什么时候、什么价、押多少钱、押
哪一边）全部抓下来，压成"每分钟市场对某件事的看法变了多少"，再检验这个
变化和中国商品期货的价格有没有关系。下面按正文出现的顺序，用大白话解释
里面的名词。已熟悉预测市场与链上数据的读者可直接跳到第 \ref{sec:acquire} 节。

\subsection{一、预测市场是什么}
\begin{description}[style=nextline, leftmargin=1em]
\item[预测市场] 一个可以拿钱押"未来某事成不成"的市场。例：押"美联储
  9 月降息"。押对的人事后按 1 元/份拿钱，押错的归零。

\item[价格就是概率] 每个结果做成一种"份额"，价格在 0 到 1 元之间。某
  份额卖 0.62 元，就等于市场认为这件事有 62\% 的概率发生；价格涨到
  0.70，就是市场把概率上调到了 70\%。这是全文的出发点：\textbf{价格 $=$
  市场此刻认为的概率}。

\item[Yes 份额 / No 份额] 一个二元问题有两种份额：押"会发生"的 Yes
  份额、押"不会发生"的 No 份额。两者价格永远加起来等于 1 元，否则有人
  能无风险套利。

\item[市场、\code{condition\_id}] 一道题就是一个"市场"，系统给它一个
  唯一编号 \code{condition\_id}，相当于这道题的身份证号。
\end{description}

\subsection{二、"链上"和"链下"}
\begin{description}[style=nextline, leftmargin=1em]
\item[区块链、Polygon] 一个公开、人人可查、不能篡改的公共账本。Polygon
  是 Polymarket 所用的那一条链。

\item[链上 / 链下] "链上" $=$ 写进了公共账本、永久可查；"链下" $=$ 只在
  公司自己服务器里、外人看不到也追不回。Polymarket 的关键设计：\textbf{谁
  挂了什么单在链下（查不到），但两张单一旦成交、钱货两清写进账本在链上
  （查得到）}。所以我们能重建每一笔成交，却永远拿不到历史挂单与盘口。

\item[订单簿、挂单、吃单] 订单簿就是一堆"我愿意用某价买/卖"的挂单排队。
  挂单方（maker，做市方）先把单子摆着等；吃单方（taker，主动方）看到
  合适的价直接成交。\textbf{吃单方是主动的一方}，这点在判断买卖方向时
  是关键。

\item[区块、区块时间戳] 链上记录按"区块"一批批打包，每个区块盖一个
  时间戳。同一区块里的所有成交共用同一个时间；实测约每 1.6 秒出一个
  区块，所以时间精度只到秒级，且同一秒内谁先谁后无法区分。

\item[撮合合约（CTF Exchange）] 链上负责"把买方卖方配对、钱货两清"的
  那段程序。它像一个自动清算中心，站在每一笔成交的中间。

\item[结算、预言机] 事件有了结果（降息了 / 没降），需要有人把"正确答案"
  写上链，这一步叫结算（\code{ConditionResolution}）；负责裁定答案的
  机制叫预言机（这里是 UMA）。结算后 Yes/No 份额价格锁死在 1 或 0。
\end{description}

\subsection{三、一笔成交里都记了什么}
\begin{description}[style=nextline, leftmargin=1em]
\item[代币、\code{asset\_id}] 每种份额（某市场的 Yes、或 No）在链上是
  一种"代币"，有唯一编号 \code{asset\_id}。买卖 Polymarket 的份额，
  本质就是买卖这些代币。

\item[\code{outcome\_seq}（第几个结果）] 标记这枚代币是这个市场的第 1 个
  结果（Yes）还是第 2 个（No）。它是代币出生就定死的属性，永不改变。

\item[USDC、名义额] USDC 是一种与美元 1:1 的数字货币，Polymarket 用它
  计价。一笔成交的名义额 $=$ 价格 $\times$ 份数（例：0.62 元 $\times$
  1000 份 $=$ 620 USDC）。

\item[\code{p\_event}（事件概率）] 把"这枚代币值多少钱"统一换算成"这件
  事发生的概率"。买的是 Yes 币，概率就是它的价；买的是 No 币，概率就是
  1 减它的价。这样不管成交发生在哪一侧，都能读成同一个"事件概率"。

\item[主动方向 \code{taker\_direction}] 这笔成交里，主动的一方（吃单方）
  是在买还是在卖。由"谁付了 USDC、谁收了代币"这个事实直接判定，不靠猜。

\item[方向 $D$（这笔在押哪边）] 把上面两件事相乘，得到"这笔主动成交是在
  押事件\emph{会}发生、还是\emph{不会}发生"：主动买 Yes 或主动卖 No $=$
  押会发生（$+1$）；主动卖 Yes 或主动买 No $=$ 押不会发生（$-1$）。
  \textbf{它只看成交本身，不看成交之后价格怎么走，不是未来信息}（详见第
  \ref{sec:dir} 节）。
\end{description}

\subsection{四、为什么有些成交要扔掉}
\begin{description}[style=nextline, leftmargin=1em]
\item[中继腿] 撮合合约站在每笔成交中间，于是账本里一笔真实成交常被记成
  两条：一条"用户对撮合合约"、一条"撮合合约对另一用户"。全加起来会把
  成交量算成两倍。判据很简单：只要一条记录的对手方是撮合合约本身，它就
  是"中继腿"，扔掉。

\item[铸造 / 合并（mint / merge）] Polymarket 允许把 1 个 USDC 拆成
  "1 份 Yes $+$ 1 份 No"（铸造），或把一对合并回 1 USDC（合并）。这是
  系统内部的拆分组合动作，没有"从别人手里接货"的对手方，不代表任何人
  看多看空。它们在账本里是\textbf{另一类记录、另一个类型标签}，本来就
  不在我们用的成交流里，谈不上剔除。

\item[negRisk] 一类特殊的多选市场（如"下任总统是谁"），机制不同，本
  数据两段都统一排除，代价是 2024 美国大选等大市场不在样本里。
\end{description}

\subsection{五、怎么把"每一笔"压成"每分钟一个数"}
\begin{description}[style=nextline, leftmargin=1em]
\item[信念创新 $N1$] 全文最核心的自造量。意思是"这一分钟里，市场对某个
  主题（如中东冲突）的看法比上一分钟变了多少"。它是一个数：正表示概率
  被推高，负表示被推低。后面几乎所有信号都由它衍生。

\item[logit 变换] 概率被卡在 0 到 1 之间，靠近两端时"从 0.98 到 0.99"
  其实是很大的信息，数值却只差 0.01。logit 是把 0 到 1 拉伸成正负无穷的
  数学变换，让两端的变化不被压扁。

\item[买卖跳动] 同一时刻成交价会在"买家愿付价"和"卖家愿收价"之间来回
  跳，这种跳动不代表概率真变了。处理办法是分钟内按买卖方向分别取中位数
  再平均，把来回跳抹平。

\item[过时报价（\code{p\_age}）] 某个冷门市场可能几小时没成交，它上一次
  的价格已经"馊了"。规则：沿用最后一次成交价，但超过 120 分钟没更新就当
  缺失，不拿陈旧价冒充新信息。

\item[时点化累计、加权] 一个主题下有很多市场，成交额大的更可信、权重更
  高。关键是权重只能用"到当前为止"累计的成交额，不能用整段时间的总额，
  否则等于让 1 月的计算偷看了 4 月才发生的事。

\item[分钟桶、右端标签] 把连续成交按每分钟归拢成一格（一个"桶"）。给这
  一格贴的时间标签取\textbf{分钟结束那一刻}，且这一格\textbf{只装该时刻
  之前的成交}。这样信号永远早于它要解释的价格，杜绝偷看未来。

\item[时序归一化（z-score）] 同一个量在不同品种、不同时期数值大小不可比。
  z-score 的做法是：减掉最近一段时间的平均值，再除以这段时间的波动大小，
  换算成"比平常高 / 低几个身位"。"最近一段"只回看过去、不含未来。

\item[稳健阈值、MAD] 判断"这一分钟的变化算不算一次事件"要有一条线。这
  条线用 MAD（一种抗极端值干扰的波动度量）乘系数得到，比用普通标准差更
  不容易被个别异常值带偏。
\end{description}

\subsection{六、中国期货那边的行话}
\begin{description}[style=nextline, leftmargin=1em]
\item[主力连续、换月] 同一品种（如原油）同时挂着好几个到期月份的合约，
  交易最活跃的那个叫"主力"。主力会随时间换到下一个月份，叫"换月"。把
  历次主力接成一条不断的价格序列，就是"主力连续"。

\item[日盘 / 夜盘] 中国商品期货一天分两段交易：白天的日盘、晚上的夜盘。
  夜盘常跨过午夜，因此\textbf{某天晚上的夜盘算作下一个交易日}。

\item[涨跌停（\code{locked}）] 一天价格最多涨 / 跌一个固定幅度，到顶就
  "封板"，全天价格钉住不动。本数据没有涨跌停价表，只能用"这一分钟最高价
  $=$ 最低价"近似判断是否封板，且只用于监控、不参与计算。

\item[连续分钟段] 期货有午休、日夜盘之间的休市。规则：只要相邻两根 K 线
  时间差不是正好 1 分钟，就断成两段。任何"未来收益"都不许跨过一个断点，
  避免把一夜或一个午休的跳空错算成 1 分钟的涨跌。

\item[前向收益 \code{fwd}] "信号发出后，价格在接下来 1、2、3、5、10、
  15 分钟涨跌了多少"，这是我们要预测的目标。

\item[基点 bp] 万分之一。10\,bp $=$ 千分之一 $=$ 0.1\%。期货分钟级涨跌
  很小，用 bp 计更清楚。
\end{description}

\subsection{七、怎么判断一个信号到底有没有用}
\begin{description}[style=nextline, leftmargin=1em]
\item[防前视、未来函数] "未来函数"指计算时不小心用到了当时还不知道的
  信息，会造出虚假的准确度、实盘一分钱赚不到。全文反复强调的"防前视"
  就是层层堵死这种偷看。

\item[RankIC（秩相关）] 衡量"信号大的时候、之后的收益是不是也偏大"的
  相关系数，范围 $-1$ 到 $1$，0 表示没关系。用"排名"而非原值计算，抗
  极端值。分钟级上 0.01 到 0.03 就已经算有一点点关系。

\item[ICIR] 把"每天的 RankIC"都求出来，再看它\textbf{是否稳定为正}：
  平均值除以波动。数值高说明信号不是靠某几天的运气，而是天天都灵。

\item[事件研究、基准] 对"发出 / 没发出"的离散信号，看它触发之后价格
  平均怎么走，并和"这个品种平时的平均走势（基准）"比。只有明显好过基准
  才算有信息。

\item[多重检验、Bonferroni 门槛] 我们试了上千个"信号 $\times$ 视界"的
  组合。试得越多，纯靠运气蒙出一个"看着显著"的就越容易。Bonferroni 是
  一条按试验次数抬高的及格线（本文约 $|t|>4$），只有跨过它才当真。
  \textbf{这正是全文最关键结论所在}：只用 Polymarket 数据的信号，没有一个
  跨过这条线。

\item[聚类稳健 $t$] 信号触发在时间上常扎堆（一件大事连着触发很多次），
  普通统计会把这些当成互相独立的证据、高估显著性。"按交易日聚类"的 $t$
  值把同一天的成簇效应扣掉，是更保守也更诚实的算法。
\end{description}
